% Môn: Toán
% Lớp: 11
% Chương: Chương I. Hàm số lượng giác và phương trình lượng giác
% Bài: Bài 2. Công thức lượng giác
% Số câu: 25
% App đọc file này trực tiếp — sửa rồi Commit trên GitHub, bấm Đồng bộ GitHub .tex

% Bài 2. Công thức lượng giác

% ===== Câu 1 =====
% ID: T11C1B2-01-DS
% Mức: VD
\dangbt{Xác định tính khả nghiệm của phương trình lượng giác}
\begin{ex}
% ID: T11C1B2-01-DS
% Mức: VD
Từ một vị trí ban đầu trong không gian, vệ tinh $X$ chuyển động theo quỹ đạo là một đường tròn quanh Trái Đất và luôn cách tâm Trái Đất một khoảng bằng $9200$ km. Sau 2 giờ thì vệ tinh $X$ hoàn thành hết một vòng di chuyển. Các khẳng định sau đây đúng hay sai?
\choiceTF
{\True Quãng đường vệ tinh $X$ chuyển động được sau $1$ giờ là $\approx 28902,65$ (km)}
{\True Quãng đường vệ tinh $X$ chuyển động được sau $1,5$ giờ là $\approx 43353,98$ (km)}
{Sau khoảng 5,3 giờ thì $X$ di chuyển được quãng đường $240000$ (km)}
{Giả sử vệ tinh di chuyển theo chiều dương của đường tròn, sau $4,5$ giờ thì vệ tinh vẽ nên một góc $\dfrac{9\pi}{2}$ rad}
\loigiai{
\begin{itemchoice}
\item (Đúng) Quãng đường vệ tinh $X$ chuyển động được sau $1$ giờ là $\approx 28902,65$ (km). (Vì): Sau $1$ giờ, vệ tinh di chuyển được $\dfrac{1}{2}$ vòng, quãng đường là $\dfrac{1}{2}C = 9200\pi \approx 28902,65$ (km
\item (Đúng) Quãng đường vệ tinh $X$ chuyển động được sau $1,5$ giờ là $\approx 43353,98$ (km). (Vì): Sau $1,5$ giờ, vệ tinh di chuyển được $\dfrac{3}{4}$ vòng, quãng đường là $\dfrac{3}{4}C = \dfrac{3}{4} \cdot 18400\pi = 13800\pi \approx 43353,98$ (km
\item (Sai) Sau khoảng 5,3 giờ thì $X$ di chuyển được quãng đường $240000$ (km). (Vì): ố giờ để vệ tinh đi được $240000$ km là $\dfrac{240000}{9200\pi} \approx 8,3$ (giờ), không phải $5,3$ giờ
\item (Sai) Giả sử vệ tinh di chuyển theo chiều dương của đường tròn, sau $4,5$ giờ thì vệ tinh vẽ nên một góc $\dfrac{9\pi}{2}$ rad.
\end{itemchoice}
}
\end{ex}

% ===== Câu 2 =====
% ID: T11C1B2-01-TLN
% Mức: NB
\dangbt{Tìm giá trị lớn nhất của biểu thức lượng giác}
\begin{ex}
% ID: T11C1B2-01-TLN
% Mức: NB
Cho biểu thức $f(x)=3\left(\sin ^4 x+\cos ^4 x\right)-2\left(\sin ^6 x+\cos ^6 x\right)$. Tính $f(1)$.
\shortans{$1$}
\loigiai{
Ta có $\sin ^4 x+\cos ^4 x=1-2\sin ^2 x\cos ^2 x$, $\sin ^6 x+\cos ^6 x=1-3\sin ^2 x\cos ^2 x$.
Suy ra $f(x)=3\left(1-2\sin ^2 x\cos ^2 x\right)-2\left(1-3\sin ^2 x\cos ^2 x\right)=1$.
Vậy biểu thức $f(1) =1$.
}
\end{ex}

% ===== Câu 3 =====
% ID: T11C1B2-01-TN
% Mức: NB
\begin{ex}
% ID: T11C1B2-01-TN
% Mức: NB
Tìm khẳng định { $\bf$ sai} (với điều kiện các hệ thức đã xác định)?
\choice
{$\tan\left(\pi+\alpha\right)=\tan\alpha$}
{\True $\cos\left(\dfrac{\pi}{2}+\alpha\right)=\sin\alpha$}
{$\cot\left(-\alpha\right)=-\cot\alpha$}
{$\sin\left(\pi-\alpha\right)=\sin\alpha$}
\loigiai{
Ta có $\cos\left(\dfrac{\pi}{2}+\alpha\right)=-\sin\alpha$.
}
\end{ex}

% ===== Câu 4 =====
% ID: T11C1B2-02-DS
% Mức: VD
\dangbt{Xác định tính khả nghiệm của phương trình lượng giác}
\begin{ex}
% ID: T11C1B2-02-DS
% Mức: VD
Các khẳng định sau đây đúng hay sai?
\choiceTF
{\True $A=\dfrac{1}{\sqrt{3}}\cos{30^{^\circ}}-3\sqrt{2}\sin{45^{^\circ}}+\cot{45^{^\circ}}=-\dfrac{3}{2}$}
{\True $B=\sin^260^{^\circ}+\tan^230^{^\circ}-2=-\dfrac{11}{12}$}
{\True $C=2\sin\dfrac{\pi}{6}+\sqrt{3}\cos\dfrac{\pi}{6}=\dfrac{5}{2}$}
{$D=\dfrac{\tan\dfrac{\pi}{4}+1}{\dfrac{1}{\sqrt{2}}\cot\dfrac{\pi}{2}-2}=\dfrac{6}{5}$}
\loigiai{
\begin{itemchoice}
\item (Đúng) $A=\dfrac{1}{\sqrt{3}}\cos{30^{^\circ}}-3\sqrt{2}\sin{45^{^\circ}}+\cot{45^{^\circ}}=-\dfrac{3}{2}$. (Vì): Ta có $A=\dfrac{1}{\sqrt{3}}\cos{30^{^\circ}}-3\sqrt{2}\sin{45^{^\circ}}+\cot{45^{^\circ}} = \dfrac{1}{\sqrt{3}}\cdot\dfrac{\sqrt{3}}{2}-3\sqrt{2}\cdot\dfrac{\sqrt{2}}{2}+1 = \dfrac{1}{2}-3+1 = -\dfrac{3}{2}$
\item (Đúng) $B=\sin^260^{^\circ}+\tan^230^{^\circ}-2=-\dfrac{11}{12}$. (Vì): Ta có $B=\sin^260^{^\circ}+\tan^230^{^\circ}-2 = \left(\dfrac{\sqrt{3}}{2}\right)^2+\left(\dfrac{1}{\sqrt{3}}\right)^2-2 = \dfrac{3}{4}+\dfrac{1}{3}-2 = \dfrac{9+4-24}{12} = -\dfrac{11}{12}$
\item (Đúng) $C=2\sin\dfrac{\pi}{6}+\sqrt{3}\cos\dfrac{\pi}{6}=\dfrac{5}{2}$. (Vì): Ta có $C=2\sin\dfrac{\pi}{6}+\sqrt{3}\cos\dfrac{\pi}{6} = 2\cdot\dfrac{1}{2}+\sqrt{3}\cdot\dfrac{\sqrt{3}}{2} = 1+\dfrac{3}{2} = \dfrac{5}{2}$
\item (Sai) $D=\dfrac{\tan\dfrac{\pi}{4}+1}{\dfrac{1}{\sqrt{2}}\cot\dfrac{\pi}{2}-2}=\dfrac{6}{5}$. (Vì): Ta có $D=\dfrac{\tan\dfrac{\pi}{4}+1}{\dfrac{1}{\sqrt{2}}\cot\dfrac{\pi}{2}-2} = \dfrac{1+1}{\dfrac{1}{\sqrt{2}}\cdot 0-2} = \dfrac{2}{-2} = -1$. Giá trị này khác với $-\dfrac{6}{5}$
\end{itemchoice}
}
\end{ex}

% ===== Câu 5 =====
% ID: T11C1B2-02-TLN
% Mức: VD
\dangbt{Tính giá trị cot từ giá trị tan đã cho}
\begin{ex}
% ID: T11C1B2-02-TLN
% Mức: VD
Cho $\tan\alpha+\cot\alpha=m$. Tìm $m>0$ để $\tan ^2\alpha+\cot ^2\alpha=7$.
\shortans{$3$}
\loigiai{
Ta có $7=\tan ^2\alpha+\cot ^2\alpha=(\tan\alpha+\cot\alpha)^2-2\tan\alpha\cot\alpha=m^2-2$.
Suy ra $m^2=9\Rightarrow m=3$.
}
\end{ex}

% ===== Câu 6 =====
% ID: T11C1B2-02-TN
% Mức: TH
\dangbt{Tìm giá trị lớn nhất của biểu thức lượng giác}
\begin{ex}
% ID: T11C1B2-02-TN
% Mức: TH
Tìm khẳng định { $\bf$ sai} trong các khẳng định sau đây?
\choice
{$\tan 45^\circ \lt \tan 60^\circ$}
{$\cos 45^\circ\le\sin 45^\circ$}
{$\sin 60^\circ \lt \sin 80^\circ$}
{\True $\cos 35^\circ \gt \cos 10^\circ$}
\loigiai{
Khi $\alpha\in\left(0^\circ\,;\,90^\circ\right)$ hàm $\cos\alpha$ là hàm giảm nên $\cos 35^\circ \lt \cos 10^\circ$ suy ra $\cos 35^\circ \gt \cos 10^\circ$ sai.
}
\end{ex}

% ===== Câu 7 =====
% ID: T11C1B2-03-TLN
% Mức: VD
\dangbt{Tính giá trị cot từ giá trị tan đã cho}
\begin{ex}
% ID: T11C1B2-03-TLN
% Mức: VD
Cho $\cos\alpha=\dfrac{3}{4}$. Tính giá trị của biểu thức $B=\dfrac{\tan\alpha+3\cot\alpha}{\tan\alpha+\cot\alpha}$ (làm tròn đến phần trăm).
\shortans{$2{,}13$}
\loigiai{
Ta có $\cos\alpha=\dfrac{3}{4}\Rightarrow 1+\tan ^2\alpha=\dfrac{1}{\cos ^2\alpha}=\dfrac{16}{9}\Rightarrow\tan ^2\alpha=\dfrac{7}{9}$.
Khi đó $B=\dfrac{\tan\alpha+\dfrac{3}{\tan\alpha}}{\tan\alpha+\dfrac{1}{\tan\alpha}}=\dfrac{\tan ^2\alpha+3}{\tan ^2\alpha+1}=\dfrac{\dfrac{7}{9}+3}{\dfrac{7}{9}+1}=\dfrac{17}{8}=2,13$.
}
\end{ex}

% ===== Câu 8 =====
% ID: T11C1B2-03-TN
% Mức: VD
\dangbt{Tìm giá trị lớn nhất của biểu thức lượng giác}
\begin{ex}
% ID: T11C1B2-03-TN
% Mức: VD
Biết $\cos\alpha=-\dfrac{1}{3}$. Khi đó, giá trị biểu thức $A=\dfrac{5\cot\alpha+4\tan\alpha}{5\cot\alpha-4\tan\alpha}$ là
\choice
{\True $-\dfrac{37}{27}$}
{$\dfrac{37}{27}$}
{$\dfrac{15}{27}$}
{$\dfrac{19}{27}$}
\loigiai{
Ta có $1+\tan^2\alpha=\dfrac{1}{\cos^2\alpha}\Rightarrow\tan^2\alpha=\dfrac{1}{\cos^2\alpha}-1=\dfrac{1}{\left(-\dfrac{1}{3}\right)^2}-1=8$.
 $A=\dfrac{5\cot\alpha+4\tan\alpha}{5\cot\alpha-4\tan\alpha}=\dfrac{5+4\tan^2\alpha}{5-4\tan^2\alpha}=\dfrac{5+4\cdot8}{5-4\cdot8}=-\dfrac{37}{27}$.
}
\end{ex}

% ===== Câu 9 =====
% ID: T11C1B2-04-TLN
% Mức: VD
\dangbt{Xác định tính đúng sai của các hệ thức lượng giác từ tổng sin cos}
\begin{ex}
% ID: T11C1B2-04-TLN
% Mức: VD
Biểu thức sau $T=2\sin\left(\dfrac{9\pi}{2}-x\right)+3\cos (19\pi-x)=k\cos x$. Khi đó $k$ bằng bao nhiêu?
\shortans{$-1$}
\loigiai{
Ta có 
$\begin{aligned}
 T&=2\sin\left(4\pi+\dfrac{\pi}{2}-x\right)+3\cos (18\pi+\pi-x)\
 &=2\sin\left(\dfrac{\pi}{2}-x\right)+3\cos (\pi-x)=2\cos x-3\cos x=-\cos x \Rightarrow k=-1.
 \end{aligned}$
}
\end{ex}

% ===== Câu 10 =====
% ID: T11C1B2-04-TN
% Mức: VD
\dangbt{Tìm giá trị lớn nhất của biểu thức lượng giác}
\begin{ex}
% ID: T11C1B2-04-TN
% Mức: VD
Giá trị lớn nhất của biểu thức $T=\cos^2x-2\sin^2x$ là
\choice
{$-2$}
{$4$}
{\True $1$}
{$2$}
\loigiai{
$T=1-\sin^2x-2\sin^2x=1-3\sin^2x$.\
  Ta có 0\le{\sin^2}x\le 1\Leftrightarrow 0 $\ge-3\sin^2x\ge-3,\forall$ x $\inR$ \Leftrightarrow 1\ge 1-3\sin^2x\ge-2 $.
Giá trị lớn nhất là 1.$
}
\end{ex}

% ===== Câu 11 =====
% ID: T11C1B2-05-TLN
% Mức: VD
\dangbt{Xác định tính đúng sai của các hệ thức lượng giác từ tổng sin cos}
\begin{ex}
% ID: T11C1B2-05-TLN
% Mức: VD
Cho hai góc nhọn $a$ và $b$. Biết $\cos a=\dfrac{1}{3}$ và $\cos b=\dfrac{1}{4}$. Tính giá trị của
 $P=(\cos a\cdot\cos b)^2-(\sin a\cdot\sin b)^2 .$ (làm tròn đến chục)
\shortans{$-0{,}8$}
\loigiai{
Ta có
$\begin{aligned}
  {P}&=(\cos a\cos b)^2-(\sin a\sin b)^2=(\cos a\cos b)^2-\left(1-\cos ^2 a\right)\left(1-\cos ^2 b\right)\
  &=\left(\dfrac{1}{12}\right)^2-\dfrac{8}{9}\cdot\dfrac{15}{16}=-\dfrac{119}{144}\approx-0,8.
  \end{aligned}$
}
\end{ex}

% ===== Câu 12 =====
% ID: T11C1B2-05-TN
% Mức: VD
\dangbt{Tìm giá trị lớn nhất của biểu thức lượng giác}
\begin{ex}
% ID: T11C1B2-05-TN
% Mức: VD
Giá trị của biểu thức $S=3-\sin^290^\circ+2\cos^260^\circ-3\tan^245^\circ$ bằng
\choice
{$\dfrac{1}{2}$}
{\True $-\dfrac{1}{2}$}
{$1$}
{$3$}
\loigiai{
Ta có $S=3-\sin^290^\circ+2\cos^260^\circ-3\tan^245^\circ =3-1^2+2\cdot \left(\dfrac{1}{2}\right)^2-3\cdot 1^2=-\dfrac{1}{2}$.
}
\end{ex}

% ===== Câu 13 =====
% ID: T11C1B2-06-TLN
% Mức: VD
\dangbt{Xác định tính đúng sai của các hệ thức lượng giác từ tổng sin cos}
\begin{ex}
% ID: T11C1B2-06-TLN
% Mức: VD
Biểu thức $A=\tan\left(\dfrac{17\pi}{2}-x\right)+2\cot (5\pi+x)=k\cot x$, Khi đó $k$ bằng bao nhiêu?
\shortans{$3$}
\loigiai{
$A=\tan\left(8\pi+\dfrac{\pi}{2}-x\right)+2\cot x=\tan\left(\dfrac{\pi}{2}-x\right)+2\cot x=\cot x+2\cot x=3\cot x$.
$\Rightarrow k=3.$
}
\end{ex}

% ===== Câu 14 =====
% ID: T11C1B2-06-TN
% Mức: VDC
\dangbt{Tính tổng các bình phương của sin với các góc cách đều}
\begin{ex}
% ID: T11C1B2-06-TN
% Mức: VDC
Tính $S=\sin^25^0+\sin^210^0+\sin^215^0+ \cdots \sin^28^0+\sin^285^0$
\choice
{$\dfrac{19}{2}$}
{$8$}
{\True $\dfrac{17}{2}$}
{$9$}
\loigiai{
$\begin{aligned}
  S&=\sin^25^0+\sin^210^0+\sin^215^0+\cdots +\sin^280^0+\sin^285^0\
  &=\left(\sin^25^0+\sin^285^0\right)+\left(\sin^210^0+\sin^280^0\right)+ \cdots + \left(\sin^215^0+\sin^275^0\right)+\sin^245^0\
  &=8+\dfrac{1}{2}=\dfrac{17}{2}.
  \end{aligned}$
}
\end{ex}

% ===== Câu 15 =====
% ID: T11C1B2-07-TLN
% Mức: VD
\dangbt{Xác định tính đúng sai của các hệ thức lượng giác từ tổng sin cos}
\begin{ex}
% ID: T11C1B2-07-TLN
% Mức: VD
Biết $\sin a+\cos a=\sqrt{2}$. Tính giá trị của $\sin ^4 a+\cos ^4 a$.
\shortans{$0{,}5$}
\loigiai{
Ta có $\sin a+\cos a =\sqrt{2}\Rightarrow 2=(\sin a+\cos a)^2 = \sin ^2 a+2\sin a\cos a+\cos ^2 a =1+2\sin a\cos a$, suy ra $\sin a\cos a=\dfrac{1}{2}$.
Khi đó $\sin ^4 a+\cos ^4 a=(\sin ^2 a+\cos ^2 a)^2-2\sin ^2 a\cos ^2 a=1-2\left(\dfrac{1}{2}\right)^2=\dfrac{1}{2}=0,5$.
}
\end{ex}

% ===== Câu 16 =====
% ID: T11C1B2-07-TN
% Mức: TH
\dangbt{Xác định giá trị sin, cos từ cot đã cho và khoảng góc}
\begin{ex}
% ID: T11C1B2-07-TN
% Mức: TH
Cho $\dfrac{\pi}{2}\lt \alpha \lt \pi$. Kết quả đúng là
\choice
{$\sin\alpha \gt 0,\cos\alpha \gt 0$}
{$\sin\alpha \lt 0,\cos\alpha \lt 0$}
{\True $\sin\alpha \gt 0,\cos\alpha \lt 0$}
{$\sin\alpha \lt 0,\cos\alpha \gt 0$}
\loigiai{
Vì $\dfrac{\pi}{2}\lt \alpha \lt \pi$ nên $\alpha$ nằm trong góc phần tư thứ $2$ nên $\sin\alpha \gt 0,\cos\alpha \lt 0$.
}
\end{ex}

% ===== Câu 17 =====
% ID: T11C1B2-08-TLN
% Mức: VD
\dangbt{Xác định tính đúng sai của các hệ thức lượng giác từ tổng sin cos}
\begin{ex}
% ID: T11C1B2-08-TLN
% Mức: VD
Cho $\cot\alpha=\dfrac{1}{3}$. Tính giá trị của biểu thức $A=\dfrac{3\sin\alpha+4\cos\alpha}{2\sin\alpha-5\cos\alpha}$.
\shortans{13}
\loigiai{
Vì $\cot\alpha=\dfrac{1}{3}$ nên $\sin\alpha\neq 0$. Chia cả tử và mẫu của biểu thức $A$ cho $\sin\alpha$, ta có
$A=\dfrac{3+4\cot\alpha}{2-5\cot\alpha}=\dfrac{3+4\cdot\dfrac{1}{3}}{2-5\cdot\dfrac{1}{3}}=13$.
}
\end{ex}

% ===== Câu 18 =====
% ID: T11C1B2-08-TN
% Mức: TH
\dangbt{Xác định giá trị sin, cos từ cot đã cho và khoảng góc}
\begin{ex}
% ID: T11C1B2-08-TN
% Mức: TH
Cho góc $\alpha$ thỏa mãn $2\pi \lt \alpha \lt \dfrac{5\pi}{2}$. Khẳng định nào sau đây { $\bf$ sai}?
\choice
{\True $\tan\alpha \lt 0$}
{$\cot\alpha \gt 0$}
{$\sin\alpha \gt 0$}
{$\cos\alpha \gt 0$}
\loigiai{
Với $2\pi \lt  \alpha \lt  \dfrac{5\pi}{2}$, ta có $\sin\alpha \gt  0$, $\cos\alpha \gt  0$, $\tan\alpha \gt  0$, $\cot\alpha \gt  0$.
}
\end{ex}

% ===== Câu 19 =====
% ID: T11C1B2-09-TLN
% Mức: VD
\dangbt{Xác định tính đúng sai của các hệ thức lượng giác từ tổng sin cos}
\begin{ex}
% ID: T11C1B2-09-TLN
% Mức: VD
Cho $\cos x=\dfrac{1}{2}$ . Tính giá trị biểu thức $P=3\sin ^2 x+4\cos ^2 x$.
\shortans{$3{,}25$}
\loigiai{
Ta có $\cos x=\dfrac{1}{2}\Rightarrow\sin ^2 x=1-\cos ^2 x=1-\dfrac{1}{4}=\dfrac{3}{4}$.
Khi đó $P=3\sin ^2 x+4\cos ^2 x=3\cdot\dfrac{3}{4}+4\cdot\dfrac{1}{4}=\dfrac{13}{4}=3,25$.
}
\end{ex}

% ===== Câu 20 =====
% ID: T11C1B2-09-TN
% Mức: TH
\dangbt{Xác định giá trị sin, cos từ cot đã cho và khoảng góc}
\begin{ex}
% ID: T11C1B2-09-TN
% Mức: TH
Cho $0\lt \alpha \lt \dfrac{\pi}{2}$. Mệnh đề nào sau đây sai?
\choice
{$\sin (\alpha\,+\,\pi)\lt 0$}
{$\cos (\alpha\,+\,\pi)\lt 0$}
{\True $\tan (\pi\,-\,\alpha)\gt 0$}
{$\cot (\pi\,-\,\alpha)\lt 0$}
\loigiai{
$\tan (\pi-\alpha)=-\tan\alpha \lt 0\,\forall\,\alpha\in\left(0\,;\,\dfrac{\pi}{2}\right)$.
}
\end{ex}

% ===== Câu 21 =====
% ID: T11C1B2-10-TLN
% Mức: VDC
\dangbt{Xác định tính đúng sai của các hệ thức lượng giác từ tổng sin cos}
\begin{ex}
% ID: T11C1B2-10-TLN
% Mức: VDC
Cho $\tan x=-\dfrac{4}{3}$ và $\dfrac{\pi}{2} \lt  x \lt  \pi$. 

Tính giá trị của biểu thức $M=\dfrac{\sin ^2 x-\cos x}{\sin x-\cos ^2 x}$ (làm tròn đến phần trăm).
\shortans{$2{,}82$}
\loigiai{
Ta có $\tan x = -\dfrac{4}{3} \Rightarrow \cos^2 x = \dfrac{1}{1+\tan^2 x} = \dfrac{1}{1+\left(-\dfrac{4}{3}\right)^2} = \dfrac{1}{1+\dfrac{16}{9}} = \dfrac{1}{\dfrac{25}{9}} = \dfrac{9}{25} \Rightarrow \cos x = \pm\dfrac{3}{5}$.

Vì $\dfrac{\pi}{2} \lt  x \lt  \pi$ nên $\cos x \lt  0$, suy ra $\cos x = -\dfrac{3}{5}$.

Khi đó $\sin x = \tan x \cdot \cos x = \left(-\dfrac{4}{3}\right)\cdot\left(-\dfrac{3}{5}\right) = \dfrac{4}{5}$.

Vậy $M = \dfrac{\sin^2 x - \cos x}{\sin x - \cos^2 x} = \dfrac{\left(\dfrac{4}{5}\right)^2 - \left(-\dfrac{3}{5}\right)}{\dfrac{4}{5} - \left(-\dfrac{3}{5}\right)^2} = \dfrac{\dfrac{16}{25} + \dfrac{3}{5}}{\dfrac{4}{5} - \dfrac{9}{25}} = \dfrac{\dfrac{16}{25} + \dfrac{15}{25}}{\dfrac{20}{25} - \dfrac{9}{25}} = \dfrac{\dfrac{31}{25}}{\dfrac{11}{25}} = \dfrac{31}{11} \approx 2,82$.
}
\end{ex}

% ===== Câu 22 =====
% ID: T11C1B2-10-TN
% Mức: VD
\dangbt{Xác định giá trị sin, cos từ cot đã cho và khoảng góc}
\begin{ex}
% ID: T11C1B2-10-TN
% Mức: VD
Cho $\sin\alpha=\dfrac{3}{5}\left(\dfrac{\pi}{2}\lt \alpha \lt \pi\right)$. Giá trị của $\cos\alpha$ bằng
\choice
{\True $\cos\alpha=-\dfrac{4}{5}$}
{$\cos\alpha=\dfrac{4}{5}$}
{$\cos\alpha=\dfrac{2}{5}$}
{$\cos\alpha=-\dfrac{2}{5}$}
\loigiai{
Ta có $\sin^2\alpha+\cos^2\alpha=1\Rightarrow{\cos^2}\alpha=1-\sin^2\alpha=1-\dfrac{9}{25}=\dfrac{16}{25}$.
Mà $\dfrac{\pi}{2}\lt \alpha \lt \pi\Rightarrow\cos\alpha=-\dfrac{4}{5}$.
}
\end{ex}

% ===== Câu 23 =====
% ID: T11C1B2-11-TLN
% Mức: VDC
\dangbt{Xác định tính đúng sai của các hệ thức lượng giác từ tổng sin cos}
\begin{ex}
% ID: T11C1B2-11-TLN
% Mức: VDC
Cho tam giác $ABC$, khi đó biểu thức $\dfrac{\sin^3 \dfrac{B}{2}}{\cos \left( \dfrac{A+C}{2} \right)} + \dfrac{\cos^3 \dfrac{B}{2}}{\sin \left( \dfrac{A+C}{2} \right)} - \dfrac{\cos (A+C)}{\sin B} \cdot \tan B$ bằng?
\shortans{2}
\loigiai{
Vì $A+B+C=180^{\circ}$ nên $A+C=180^{\circ}-B$.
$\begin{aligned}
  \text{VT}& =\dfrac{\sin^3\dfrac{B}{2}}{\cos\left(\dfrac{180^{^\circ}-B}{2}\right)}+\dfrac{\cos^3\dfrac{B}{2}}{\sin\left(\dfrac{180^{^\circ}-B}{2}\right)}-\dfrac{\cos\left(180^\circ-B\right)}{\sin B}\cdot\tan B\ 
  &=\dfrac{\sin^3\dfrac{B}{2}}{\sin\dfrac{B}{2}}+\dfrac{\cos^3\dfrac{B}{2}}{\cos\dfrac{B}{2}}-\dfrac{-\cos B}{\sin B}\cdot\tan B\ 
  &=\sin^2\dfrac{B}{2}+\cos^2\dfrac{B}{2}+\cot B\cdot\tan B=1+1=2=\text{VP}.
  \end{aligned}$
}
\end{ex}

% ===== Câu 24 =====
% ID: T11C1B2-11-TN
% Mức: NB
\dangbt{Xác định tính khả nghiệm của phương trình lượng giác}
\begin{ex}
% ID: T11C1B2-11-TN
% Mức: NB
Cho $a$ là số thực bất kỳ. Chọn khẳng định đúng trong các khẳng định sau:
\choice
{$\sin a+\cos a=1$}
{$\sin^3a+\cos^3a=1$}
{$\sin^4a+\cos^4a=1$}
{\True $\sin^2a+\cos^2a=1$}
\loigiai{
Ta có $\sin^2a+\cos^2a=1$ với mọi $a\in\mathbb{R}$ nên $\sin^2a+\cos^2a=1$ đúng.
}
\end{ex}

% ===== Câu 25 =====
% ID: T11C1B2-12-TN
% Mức: NB
\begin{ex}
% ID: T11C1B2-12-TN
% Mức: NB
Trong các hệ thức sau đây, hệ thức nào đúng?
\choice
{\True $\sin^2x+\cos^2x=1$}
{$\sin{x^2}+\cos{x^2}=1$}
{$\sin 2x+\cos 2x=1$}
{$\sin^2x+\cos{x^2}=1$}
\loigiai{
$\sin^2x+\cos^2x=1$.
}
\end{ex}
